\begin{table}[htbp]
\centering
\caption{Standardized Effect Sizes}
\label{tab:sde}
\adjustbox{max width=\textwidth}{\begin{tabular}{lcccccc}
\toprule
Outcome & $\hat{\beta}$ & SE & SD($Y$) & SDE & SE(SDE) & Classification \\
\midrule
\multicolumn{7}{l}{\textit{Panel A: Pooled}} \\
Log appraisal value & -0.0526 & 0.0192 & 0.3321 & -0.1584 & 0.0580 & Large negative \\
Appraisal value (\euro/m$^2$) & 19.8 & 26.8 & 421.5 & 0.0469 & 0.0636 & Small positive \\
\addlinespace
\multicolumn{7}{l}{\textit{Panel B: Heterogeneous (sample splits)}} \\
Log appraisal (Lisbon metro) & -0.0651 & 0.0403 & 0.3435 & -0.1896 & 0.1173 & Large negative \\
Log appraisal (Porto/Algarve) & -0.0433 & 0.0085 & 0.3111 & -0.1392 & 0.0275 & Moderate negative \\
\bottomrule
\end{tabular}}
\vspace{0.5em}
\begin{minipage}{\textwidth}\emergencystretch=3em
\footnotesize
\textit{Notes:} \textbf{Country:} Portugal. \textbf{Research question:} Does restricting golden visa real estate investment in prime urban areas reduce housing prices in restricted regions and divert capital to eligible interior regions? \textbf{Policy mechanism:} Decreto-Lei 14/2021 (effective January 2022) barred golden visa applicants from purchasing residential property in the Metropolitan Areas of Lisbon and Porto and in the Algarve, while retaining eligibility in interior and island NUTS3 regions, thereby eliminating a demand channel from approximately 1,000 annual foreign investor-visa property transactions in high-demand coastal markets. \textbf{Outcome definition:} Median bank housing appraisal value in EUR per square meter from INE Portugal's monthly BPIHE survey of bank property valuations. \textbf{Treatment:} Binary; NUTS3 regions classified as restricted if they lost golden visa real estate eligibility in January 2022. \textbf{Data:} INE BPIHE monthly bank appraisal survey, January 2015--December 2024, 99 municipalities, 11,553 municipality-month observations. \textbf{Method:} Two-way fixed effects DiD with municipality and year-month fixed effects plus municipality-specific linear time trends; standard errors clustered at NUTS3 level (27 clusters). \textbf{Sample:} All NUTS3 regions in mainland Portugal plus Azores and Madeira with at least 80\% monthly coverage in the BPIHE survey. SDE $= \hat{\beta} / \text{SD}(Y)$ where SD($Y$) is the pre-treatment standard deviation of the restricted group. Classification refers to magnitude, not statistical significance: Large ($|$SDE$| > 0.15$), Moderate ($0.05$--$0.15$), Small ($0.005$--$0.05$), Null ($< 0.005$).
\end{minipage}
\end{table}
